% page 268 of the facsimile (image p-267 of the scan)
% block 160.0 x 246.3 mm ; left margin 8.0 mm
\pgc{-12.700mm}{0.000mm}{
\l{\cel{5}260.}
\l{}
\l{\sou{kanonizar}. (\sou{trans}.)(\sou{religio katol.}). I. Enskribar la nomo}
\l{\cel{3}en la katalogo de la santi ; plasizar inter la santi.-}
\l{\cel{3}II. Deklarar konforma ye la kanoni ekleziala. - DEFIS.}
\l{}
\l{\sou{kanoto}. Mikra batelo, sen-ferdeka, movata per segli, remi-}
\l{\cel{3}li, motoro, ed ordinare uzata sur la mar-aquo. - FIS.}
\l{}
\l{\sou{kansono}. Ensemblo ek mikra versi, maxim-multa-kaze dividi-}
\l{\cel{3}ta ad kupleti kun refreno, kantata per ario popularo. FIS}
\l{}
\l{\sou{kantar}.\cel{2}(\sou{trans.}) I. Produktar per la voco muzik-ario. -II.}
\l{\cel{3}Celebrar per liriko-versi poeziala. - DEFIS.}
\l{}
\l{\sou{kantabilo}. (\sou{muziko)}.\cel{2}Kantajo ye movimento moderata.- DEFI.}
\l{}
\l{\sou{kantarelo.}\cel{2}(\sou{bot.)}\cel{2}Speco de boleto, fungo manjebla. - L.}
\l{\cel{3}\sou{cantharellus cibarius}. - DEFSL.}
\l{}
\l{\sou{kantarido}.\cel{2}(\sou{zool.}) Insekto koleoptera, quan onu sekigas e}
\l{\cel{3}reduktas a pulvero, uzata farmaciale en la preparaji vezik-}
\l{\cel{3}igiva. - L.\cel{1}\sou{cantharis vesicatoria}. DEFIRS.}
\l{}
\l{\sou{kantato.}\cel{2}Ensemblo ek versi, qua konsistas (generale) ye re-}
\l{\cel{3}citativo e ye strofi, a qua onu povas adjuntar muzik-ario.}
\l{\cel{3}- DEFIRS.}
\l{}
\l{\sou{kantiko}\cel{1}(\sou{kristanismo}). I. Kantajo religiala ek la santa skri-}
\l{\cel{3}buri, por laudar, por dankar, deo. - II. Kantajo liturgia-}
\l{\cel{3}la.\cel{2}- III. Kantiko spiritala ; kantajo religiala, en lin-}
\l{\cel{3}guo vulgara, a qua onu adaptas arii, e quan kantas la pueri}
\l{\cel{3}katekizata, la adolecantini di la kunfratari. - DEFIS.}
\l{}
\l{\sou{kantileno}. Speco di kansono, anciena e lamentala. - DEFIRS.}
\l{}
\l{\sou{kantino}. Taverno di regimento, di kazerno ; taverno di fabrik-}
\l{\cel{3}eyo, di karcero, e c. - DEFIS.}
\l{}
\l{\sou{kantono}.\cel{2}(En\cel{1}\sou{Francia}).\cel{2}Dividuro teritoriala di distrikto. -}
\l{\cel{3}II. (En\cel{1}\sou{Suisia}). Singla ek la 22 mikra stati,ye qui konsis-}
\l{\cel{3}tas la federitaro (o : federuro) Suisa. - DEFIRS.}
\l{}
\l{\sou{kantoro.}\cel{2}(\sou{kristanismo)}. La persono qua kantas lor la oficio}
\l{\cel{3}en kirko, en templo. - DEFIRS.}
\l{}
\l{\sou{kanulo}. Tubo, plu o min longa, rigida o flexebla, quan onu}
\l{\cel{3}uzas en ula operaci kirurgiala, o por enduktar (aden la}
\l{\cel{3}intestini o la vagino) ula liquidi. - DEFIS.}
\l{}
\l{\sou{kanvaso}.\cel{2}I.Dika telo kruda, de qua onu facas la segli, la}
\l{\cel{3}visho-tuki koqueyala. - II. Reto ek kotono-fili duopla,}
\l{\cel{3}qui interkrucumas ajure, formacante quadrati, ye qua kon-}
\l{\cel{3}sistas la fundo sur qua onu facas la agulo-tapeto. - DEFIRS.}
\l{}
\l{kaolino. Arjilo blanka, tre pura, quan onu renkontras preci-}
\l{\cel{3}pue en Chinia, ed ek qua onu facas la porcelano, - DEFIRS.}
}
